\documentclass[11pt]{beamer}

\usetheme{Madrid}
\usecolortheme{default}
\definecolor{unicaucaverde}{RGB}{0,90,60}
\setbeamercolor{structure}{fg=unicaucaverde}
\setbeamertemplate{navigation symbols}{}

\usepackage[utf8]{inputenc}
\usepackage[spanish]{babel}
\usepackage{amsmath,amssymb}
\usepackage{listings}
\usepackage{xcolor}

\lstdefinestyle{cstyle}{
  language=C,
  basicstyle=\ttfamily\scriptsize,
  keywordstyle=\color{unicaucaverde}\bfseries,
  commentstyle=\color{gray}\itshape,
  stringstyle=\color{orange!70!black},
  breaklines=true,
  showstringspaces=false,
  tabsize=2
}
\lstset{style=cstyle}

\title[Semana 2: Estructuras de control]{Estructuras de control: condicionales y ciclos}
\subtitle{Programación --- Semana 2 de 16}
\author{Universidad del Cauca}
\date{2026}

\begin{document}

\begin{frame}
  \titlepage
\end{frame}

\begin{frame}{Semana 2: dónde estamos}
  \begin{itemize}
    \item Continuación de Fundamentos de programación en C.
    \item Subtemas 2.1--2.2 --- RA2.1, RA2.2.
    \item 4 horas: condicionales y \texttt{switch}, luego ciclos y anidamiento.
  \end{itemize}
  \tableofcontents
\end{frame}

\section{Condicionales}

\begin{frame}[fragile]{\texttt{if}, \texttt{else if}, \texttt{else}}
\begin{lstlisting}
if (nota >= 4.5f) {
    printf("Excelente\n");
} else if (nota >= 3.0f) {
    printf("Aprobado\n");
} else {
    printf("Reprobado\n");
}
\end{lstlisting}
  \begin{itemize}
    \item Se evalúan en orden; se ejecuta el primer bloque verdadero.
    \item Operador ternario: \texttt{cond ? valor\_si\_true : valor\_si\_false}.
  \end{itemize}
\end{frame}

\begin{frame}{Condiciones compuestas}
  \begin{itemize}
    \item \texttt{\&\&} (y), \texttt{||} (o), \texttt{!} (no).
    \item Cortocircuito: \texttt{a \&\& b} no evalúa \texttt{b} si \texttt{a} ya es falso.
  \end{itemize}
  \begin{block}{Ejemplo}
    \texttt{if (edad >= 18 \&\& tiene\_carnet) \{ ... \}}
  \end{block}
  \vfill
  \textbf{Taller (hora 2):} menú con \texttt{switch-case}.
\end{frame}

\begin{frame}[fragile]{\texttt{switch-case}}
\begin{lstlisting}
switch (opcion) {
    case 1: resultado = a + b; break;
    case 2: resultado = a - b; break;
    default: printf("Invalido\n");
}
\end{lstlisting}
  \begin{alertblock}{Cuidado}
    Sin \texttt{break}, la ejecución \emph{cae} al siguiente \texttt{case}.
  \end{alertblock}
\end{frame}

\section{Ciclos}

\begin{frame}[fragile]{\texttt{while} vs. \texttt{do-while}}
  \begin{columns}
    \begin{column}{0.5\textwidth}
      \textbf{while} (condición al inicio)
\begin{lstlisting}
while (cond) {
    ...
}
\end{lstlisting}
      Puede ejecutarse 0 veces.
    \end{column}
    \begin{column}{0.5\textwidth}
      \textbf{do-while} (condición al final)
\begin{lstlisting}
do {
    ...
} while (cond);
\end{lstlisting}
      Se ejecuta al menos 1 vez.
    \end{column}
  \end{columns}
\end{frame}

\begin{frame}[fragile]{\texttt{for}: conteo explícito}
\begin{lstlisting}
for (int i = 1; i <= 5; i++) {
    printf("%d\n", i);
}
\end{lstlisting}
  Inicialización + condición + incremento, todo en un solo lugar. Ideal cuando se conoce el número
  de repeticiones.
\end{frame}

\begin{frame}[fragile]{\texttt{break} y \texttt{continue}}
\begin{lstlisting}
for (int i = 1; i <= 10; i++) {
    if (i == 7) break;         // termina el ciclo
    if (i % 2 == 0) continue;  // salta esta repeticion
    printf("%d\n", i);
}
\end{lstlisting}
  \begin{itemize}
    \item \texttt{break}: termina el ciclo (o \texttt{switch}) inmediatamente.
    \item \texttt{continue}: salta al resto del cuerpo, pasa a la siguiente repetición.
  \end{itemize}
\end{frame}

\begin{frame}[fragile]{Ciclos anidados}
\begin{lstlisting}
for (int fila = 1; fila <= 5; fila++) {
    for (int col = 1; col <= 5; col++) {
        printf("%3d", fila * col);
    }
    printf("\n");
}
\end{lstlisting}
  Por cada repetición del ciclo externo, el ciclo interno se ejecuta completo.
  \vfill
  \textbf{Taller (hora 4):} tabla de multiplicar, patrones, validación de entrada con
  \texttt{while}.
\end{frame}

\begin{frame}{Cierre de la semana}
  \begin{itemize}
    \item RA2.1: condicionales simples, encadenados, \texttt{switch}, condiciones compuestas.
    \item RA2.2: \texttt{while}/\texttt{do-while}/\texttt{for}, \texttt{break}/\texttt{continue},
          anidamiento.
    \item Control formativo al inicio de la Semana 3.
  \end{itemize}
\end{frame}

\end{document}
